\documentclass[conference]{IEEEtran}

\usepackage{amsmath}
\usepackage{amssymb}
\usepackage{graphicx}
\usepackage{booktabs}
\usepackage{hyperref}
\usepackage{xcolor}
\usepackage{microtype}
\usepackage{cite}
\usepackage{algorithm}
\usepackage{algorithmic}
\usepackage{multirow}

\hypersetup{colorlinks=true, linkcolor=blue, citecolor=blue, urlcolor=blue,}

\title{A Hybrid Residual Forecasting Framework for Time Series with MLOps Integration}

\author{ \IEEEauthorblockN{Vansh Kumar Soni} \IEEEauthorblockA{ Department of CSE (AI\&ML)\\ Dayananda Sagar University\\ Bangalore, India\\ eng23am0081@dsu.edu.in }
\and \IEEEauthorblockN{Vinayak K} \IEEEauthorblockA{ Department of CSE (AI\&ML)\\ Dayananda Sagar University\\ Bangalore, India\\ eng23am0086@dsu.edu.in }
\and \IEEEauthorblockN{Abhishek Tiwari} \IEEEauthorblockA{ Department of CSE (AI\&ML)\\ Dayananda Sagar University\\ Bangalore, India\\ eng23am0097@dsu.edu.in }
}

\begin{document}
  \maketitle

  % ── Abstract ──────────────────────────────────────────────────────────────────
  \begin{abstract}
    Production time series forecasting is rarely just a modelling problem---it equally
    demands reliable infrastructure for training, tracking, and serving predictions at
    scale. This paper introduces a Hybrid Residual Forecasting framework, validated on
    the ETTh1 electricity transformer benchmark. The core idea is a deliberate task
    decomposition: a Ridge regression base model handles the global linear signal, and
    then CatBoost and ExtraTrees are trained solely on the errors that Ridge leaves
    behind, targeting non-linear local fluctuations that a linear model cannot express.
    The final output adds the averaged residual correction back onto the linear base.
    Every stage of this pipeline is version-controlled with DVC, tracked in MLflow,
    packaged inside Docker, and exposed through a FastAPI REST endpoint. Against the
    ETTh1 test partition the hybrid achieves RMSE\,$=0.0757$, MAE\,$=0.0520$,
    SMAPE\,$=7.80\%$, and $\text{R}^{2}=0.9645$, beating every individual model we
    evaluated. We additionally document an important pitfall: MAPE becomes unreliable
    on standardised targets whose values hover near zero, and we argue SMAPE should
    be preferred in such settings. The complete system, from raw data download through
    live HTTP inference, runs end-to-end with a single \texttt{dvc repro} command.
  \end{abstract}

  \begin{IEEEkeywords}
    time series forecasting, hybrid forecasting, residual boosting, Ridge regression,
    CatBoost, ExtraTrees, LightGBM, XGBoost, MLOps, DVC, MLflow, FastAPI, Docker,
    reproducibility
  \end{IEEEkeywords}

  % ── 1. Introduction ───────────────────────────────────────────────────────────
  \section{Introduction}

  Forecasting plays a foundational role in energy management, financial planning, supply
  chain optimisation, and industrial monitoring \cite{makridakis2018m4, box2015time}.
  In practice, the difficulty is not limited to picking a model; it extends to building
  infrastructure that can train, retrain, track, and serve predictions with confidence
  over time. Academic work and industrial practice have historically treated these two
  concerns separately, and this disconnect is the motivation for our work.

  On the modelling side, decades of research have produced several well-understood
  families. ARIMA \cite{box2015time} models the linear autocorrelation structure of a
  series through differencing, autoregressive, and moving-average terms. Prophet
  \cite{taylor2018prophet} takes a decomposition approach, treating trend and
  seasonality as separable components. Gradient-boosted trees---LightGBM
  \cite{ke2017lightgbm} and XGBoost \cite{chen2016xgboost} in particular---reframe
  forecasting as tabular regression over lag and rolling-statistics features,
  often matching or exceeding sequence models on structured benchmarks
  \cite{hewamalage2021recurrent, januschowski2022forecasting}. None of these families
  consistently wins: statistical approaches shine on short histories, tree models
  leverage rich feature matrices, and competitions repeatedly show that combining
  different families produces more robust results than any single method
  \cite{makridakis2018m4, makridakis2022m5}.

  The operational half of the problem falls under Machine Learning Operations (MLOps)
  \cite{sculley2015hidden, kreuzberger2023mlops}. Sculley et al.\ famously noted that
  actual model code makes up a surprisingly small portion of a deployed ML system; the
  surrounding data pipelines, configuration layers, monitoring hooks, and serving
  infrastructure dominate. Modern tools exist to fill these gaps---DVC \cite{dvc2020}
  for data versioning and pipeline orchestration, MLflow \cite{mlflow2018} for
  experiment tracking and model registration, Docker \cite{docker2013} for
  environment-reproducible packaging, and FastAPI \cite{fastapi2018} for REST serving.
  A recent landscape survey \cite{berberi2025mlops} reinforces that production teams
  need coordinated support across all four of these dimensions, not just individual
  point tools.

  What is still rare in the published literature is a working system that weaves these
  tools together around a non-trivial multi-model forecasting architecture. This paper
  fills that gap.

  \textbf{Contributions.} Concretely, we contribute:
  \begin{itemize}
    \item A \textbf{Hybrid Residual Forecasting} architecture where Ridge regression
      captures the global linear trend and two tree models (CatBoost, ExtraTrees) are
      trained only on Ridge's errors---yielding $\text{R}^{2}=0.9645$ and
      SMAPE\,$=7.80\%$ on the ETTh1 benchmark.

    \item A feature engineering pipeline augmented with exponentially weighted moving
      averages (EWM) and lag-difference velocity features, designed to make local trend
      momentum visible to residual-correction models.

    \item A fully reproducible four-stage DVC pipeline, MLflow nested run tracking,
      Docker containerisation, and a live FastAPI endpoint, together with a
      quantitative analysis of MAPE's unreliability on standardised near-zero targets.
  \end{itemize}

  % ── 2. Related Work ───────────────────────────────────────────────────────────
  \section{Related Work}

  \subsection{Statistical Time Series Models}

  ARIMA \cite{box2015time} is the workhorse statistical baseline in univariate
  forecasting. It achieves stationarity through differencing and then models the
  remaining autocorrelation with AR and MA components, with model order selected by
  the Box-Jenkins procedure. Despite its theoretical elegance, ARIMA fits sequentially
  and cannot directly exploit an exogenous feature matrix---a practical limitation
  when additional covariates carry predictive signal.

  Prophet \cite{taylor2018prophet} approaches the same problem differently, fitting a
  piecewise linear or logistic trend alongside Fourier-series seasonality and holiday
  regressors through a Stan-based \cite{stan2017} Bayesian optimiser. It was designed
  primarily for daily business data with strong, regular seasonal patterns, and it
  handles irregular gaps and outliers well. Its additive structure, however, implicitly
  assumes the signal has meaningful amplitude at its original measurement scale;
  standardisation can suppress the seasonal structure that Prophet needs to function.

  \subsection{Gradient-Boosted Tree Forecasting}

  LightGBM \cite{ke2017lightgbm} departs from traditional depth-wise tree growth by
  growing leaves greedily and using histogram-based bin splitting, which dramatically
  reduces training time on large tabular datasets. XGBoost \cite{chen2016xgboost}
  takes a more regularisation-heavy approach---second-order gradient approximations,
  column subsampling, and shrinkage---leading to somewhat different generalisation
  characteristics. Both frame the forecasting task as ordinary supervised regression
  over a manually constructed lag-and-calendar feature matrix, and this simple framing
  has repeatedly been shown to rival or beat sequential neural architectures on
  structured benchmarks \cite{hewamalage2021recurrent}. The M5 competition
  \cite{makridakis2022m5} and subsequent post-competition analyses
  \cite{januschowski2022forecasting, lainder2022forecasting} solidified gradient
  boosting as the dominant algorithm family in practical large-scale forecasting.

  \subsection{Deep and Foundation Forecasting Models}

  A growing body of work revisits the assumption that deeper architectures
  automatically improve forecasting. Zeng et al.\ \cite{zeng2023transformers}
  found that single-layer linear models can actually beat much more complex
  Transformer-based forecasters on several long-horizon ETT benchmarks---a sobering
  result that motivates careful baseline construction before introducing architectural
  complexity. PatchTST \cite{nie2023patchtst} partially addresses this by segmenting
  the input series into patches and processing channels independently, achieving
  improvements on ETTh1 and related datasets. At an even larger scale, Chronos
  \cite{ansari2024chronos} tokenises time series values and pre-trains a probabilistic
  forecaster on a diverse corpus of datasets. Our work sits in a different part of
  this design space: we deliberately avoid deep architectures and instead study how
  well a two-stage linear-plus-trees pipeline can be wrapped in a production-ready
  MLOps system.

  \subsection{Forecast Combination}

  The theoretical case for combining forecasts goes back to Bates and Granger
  \cite{bates1969combination}, who demonstrated that a simple average of two unbiased
  forecasters generally reduces variance relative to either individual forecast. The M4
  \cite{makridakis2018m4} and M5 \cite{makridakis2022m5} competitions both confirmed
  this empirically at scale: hybrid methods that blend statistical and machine-learning
  forecasters consistently placed near the top of the leaderboard. More sophisticated
  combination strategies exist, including Wolpert's stacked generalisation
  \cite{wolpert1992stacked}, which trains a meta-learner on cross-validated base
  forecasts, and blending approaches that jointly tune gradient boosted trees with
  neural forecasters \cite{vogklis2022blending}. We intentionally use an unweighted
  average rather than stacking, prioritising simplicity and reduced overfitting risk
  at the cost of some potential accuracy gain.

  \subsection{MLOps Tooling}

  Kreuzberger et al.\ \cite{kreuzberger2023mlops} distilled the MLOps landscape into
  three foundational pillars: reproducibility, automation, and monitoring. A more recent
  review by Berberi et al.\ \cite{berberi2025mlops} catalogues the current generation
  of platforms and tools, concluding that no single tool yet covers all lifecycle
  concerns---data management, training orchestration, deployment, and runtime monitoring
  each tend to require dedicated solutions. DVC \cite{dvc2020} addresses the data and
  pipeline-reproducibility gap through content-addressed hashing integrated with Git.
  MLflow \cite{mlflow2018} handles experiment governance with a language-agnostic
  tracking API and an optional model registry. Docker \cite{docker2013} and FastAPI
  \cite{fastapi2018} then cover packaging and serving. The combination of all four
  in a single, coherently documented forecasting system is our specific contribution
  to this intersection of research areas.

  % ── 3. Dataset ──────────────────────────────────────────────────────────────
  \section{Dataset}

  Our experiments use the \textbf{ETTh1} (Electricity Transformer Temperature,
  hourly) dataset \cite{zhou2021informer}. Collected from an electricity transformer
  substation at one-hour resolution between July 2016 and June 2018, the corpus spans
  17,420 observations across eight channels: six load readings (HUFL, HULL, MUFL,
  MULL, LUFL, LULL) and oil temperature (OT). We target OT, whose training-set
  distribution has mean 13.32\,°C, standard deviation 8.57\,°C, and range $[-4.08,
  46.01]$\,°C.

  We partition the data chronologically to eliminate any look-ahead bias:
  \begin{itemize}
    \item \textbf{Train}: 70\% --- 12,194 samples (Jul 2016 -- Nov 2017)
    \item \textbf{Validation}: 10\% --- 1,742 samples (Nov 2017 -- Mar 2018)
    \item \textbf{Test}: 20\% --- 3,484 samples (Mar 2018 -- Jun 2018)
  \end{itemize}
  Respecting this temporal order is not merely good practice---it is a necessary
  condition for evaluating genuine out-of-sample generalisation in any forecasting
  study \cite{hyndman2018forecasting}.

  % ── 4. Methodology ──────────────────────────────────────────────────────────
  \section{Methodology}

  \subsection{Problem Formulation}

  Let $\{y_{t}\}_{t=1}^{T}$ be the raw OT sequence and
  $\mathbf{x}_{t} \in \mathbb{R}^{F}$ the corresponding feature vector ($F = 45$)
  computed from observations available up to and including time $t$. We target
  one-step-ahead prediction of $\hat{y}_{t+1}$.

  Before fitting any model we standardise the target:
  \begin{equation}
    \tilde{y}_{t} = \frac{y_{t} - \mu_{\text{train}}}{\sigma_{\text{train}}}
  \end{equation}
  where $\mu_{\text{train}}$ and $\sigma_{\text{train}}$ are computed on the training
  split exclusively. All subsequent modelling and evaluation operates on $\tilde{y}$.

  \subsection{Feature Engineering}

  The 45-dimensional feature vector $\mathbf{x}_{t}$ is built from three complementary
  families:

  \textbf{Lag and velocity features.} Direct autoregressive observations at lags
  $\mathcal{L}= \{1, 2, 3, 6, 12, 24, 48, 168\}$ hours, plus first-order differences
  between consecutive lag pairs (e.g., $\tilde{y}_{t-1} - \tilde{y}_{t-2}$). The
  difference features encode local trend direction, giving tree models an explicit
  velocity signal rather than forcing them to compute it implicitly.

  \textbf{Rolling and EWM statistics.} Window-based means over
  $\mathcal{W}= \{6, 12, 24, 168\}$ hours, together with exponentially weighted
  moving averages over the same spans. EWM places heavier weight on recent observations,
  so it reacts to sudden shifts faster than an equal-weight rolling mean.

  \textbf{Calendar encodings.} Hour of day, day of week, month of year, and a
  binary weekend indicator, each projected onto a unit circle via sine/cosine
  transforms to preserve cyclical adjacency (so that hour 23 and hour 0 are
  numerically close, not maximally distant).

  \subsection{Hybrid Architecture}

  Rather than training every model on the same target and averaging their outputs,
  we use a sequential two-phase design inspired by gradient boosting's residual-fitting
  insight.

  \subsubsection{Phase 1: Ridge Base Model}

  Tree ensembles cannot extrapolate beyond the range of their training labels, so any
  global drift in the series will systematically fool them. Ridge Regression, by
  contrast, fits a globally linear projection and extrapolates naturally. We train Ridge
  on $\tilde{\mathbf{y}}$ by minimising:
  \begin{equation}
    \mathcal{L}_{\text{Ridge}}= \|\tilde{\mathbf{y}} - \mathbf{X}\mathbf{w}\|_{2}^{2} + \alpha \|\mathbf{w}\|_{2}^{2}
  \end{equation}
  This yields a globally stable base trajectory $\hat{y}^{\text{base}}_{t}$ that
  captures the dominant linear structure of the signal.

  \subsubsection{Phase 2: Residual Correction (CatBoost \& ExtraTrees)}

  Ridge's linear constraint means it cannot capture conditional non-linearities---for
  instance, the interaction between a sustained high 24-hour rolling average and a
  specific time-of-day. We expose and target these missed patterns by computing the
  signed residual on the training set:
  \begin{equation}
    r_{t} = \tilde{y}_{t} - \hat{y}^{\text{base}}_{t}
  \end{equation}
  Both CatBoost and ExtraTrees are then fitted with $r_{t}$ as their target rather
  than the original $\tilde{y}_{t}$. CatBoost brings symmetric, depth-limited trees
  built via ordered boosting; ExtraTrees uses random threshold selection to build
  high-variance, low-bias corrective functions. Their algorithmic differences mean
  their residual predictions are partially uncorrelated, which is precisely what
  makes averaging them beneficial.

  \subsection{Hybrid Inference}

  At test time, the final prediction is:
  \begin{equation}
    \hat{y}_{t}^{\text{hybrid}}= \hat{y}_{t}^{\text{base}}+ \frac{1}{2}\left( \hat{r}_{t}^{\text{CatBoost}}+ \hat{r}_{t}^{\text{ExtraTrees}}\right)
  \end{equation}

  Algorithm~\ref{alg:hybrid} summarises the complete training and inference procedure.

  \begin{algorithm}
    \caption{Hybrid Residual Training and Inference}
    \label{alg:hybrid}
    \begin{algorithmic}
      [1] \REQUIRE Train/val/test feature matrices; targets $\tilde{\mathbf{y}}$
      \STATE \textbf{Phase 1: Base Trend}
      \STATE Fit Ridge on train features targeting $\tilde{\mathbf{y}}$
      \STATE $\hat{\mathbf{y}}^{\text{base}} \leftarrow \text{Ridge.predict}(\text{train features})$
      \STATE $\mathbf{r} \leftarrow \tilde{\mathbf{y}} - \hat{\mathbf{y}}^{\text{base}}$
      \STATE \textbf{Phase 2: Residuals}
      \STATE Fit CatBoost on train features targeting $\mathbf{r}$
      \STATE Fit ExtraTrees on train features targeting $\mathbf{r}$
      \STATE \textbf{Phase 3: Inference}
      \STATE $\hat{\mathbf{y}}^{\text{test\_base}} \leftarrow \text{Ridge.predict}(\text{test features})$
      \STATE $\hat{\mathbf{r}}^{\text{test\_cat}} \leftarrow \text{CatBoost.predict}(\text{test features})$
      \STATE $\hat{\mathbf{r}}^{\text{test\_et}} \leftarrow \text{ExtraTrees.predict}(\text{test features})$
      \STATE $\hat{\mathbf{y}}^{\text{hybrid}} \leftarrow \hat{\mathbf{y}}^{\text{test\_base}} + \frac{1}{2}(\hat{\mathbf{r}}^{\text{test\_cat}} + \hat{\mathbf{r}}^{\text{test\_et}})$
      \RETURN $\hat{\mathbf{y}}^{\text{hybrid}}$
    \end{algorithmic}
  \end{algorithm}

  \subsection{Evaluation Metrics}

  We report the following metrics on the held-out test partition. Let $y_{i}$ and
  $\hat{y}_{i}$ denote true and predicted values for $i = 1, \ldots, N$.

  \begin{equation}
    \text{MAE}= \frac{1}{N}\sum_{i=1}^{N}|y_{i} - \hat{y}_{i}|
  \end{equation}
  \begin{equation}
    \text{RMSE}= \sqrt{\frac{1}{N}\sum_{i=1}^{N}(y_{i} - \hat{y}_{i})^{2}}
  \end{equation}
  \begin{equation}
    R^{2} = 1 - \frac{\sum_{i}(y_{i} - \hat{y}_{i})^{2}}{\sum_{i}(y_{i} - \bar{y})^{2}}
  \end{equation}
  \begin{equation}
    \text{MAPE}= \frac{100}{N}\sum_{i=1}^{N}\frac{|y_{i} - \hat{y}_{i}|}{|y_{i}| +
    \varepsilon}
  \end{equation}
  \begin{equation}
    \text{SMAPE}= \frac{100}{N}\sum_{i=1}^{N}\frac{|y_{i} - \hat{y}_{i}|}{(|y_{i}|
    + |\hat{y}_{i}|)/2 + \varepsilon}
  \end{equation}

  MAPE is widely reported in the forecasting literature and is easy to explain to
  domain experts. On standardised data, however, $|y_{i}|$ is frequently close to zero,
  inflating individual terms without bound. SMAPE guards against this by symmetrising
  the denominator across both the actual and predicted magnitude, keeping the metric
  bounded in $[0\%, 200\%]$ regardless of how small the true value is.

  % ── 5. MLOps Pipeline ─────────────────────────────────────────────────────────
  \section{MLOps Pipeline}

  \subsection{Pipeline Structure and DVC}

  The pipeline is broken into four sequential DVC stages:
  \[
    \texttt{ingest}\to \texttt{preprocess}\to \texttt{featurize}\to \texttt{train}
  \]
  Every stage is declared in \texttt{dvc.yaml} with its upstream data dependencies,
  tracked \texttt{params.yaml} keys, and declared outputs. DVC hashes all inputs at
  execution time and records the result in \texttt{dvc.lock}; on the next run it skips
  any stage whose hash is unchanged. The practical benefit is fast iterative
  experimentation: updating a lag period in \texttt{params.yaml} re-runs only
  \texttt{featurize} and \texttt{train}. Committing \texttt{dvc.lock} to Git
  simultaneously snapshots the code, the parameters, and the resulting data state:
  \[
    \texttt{dvc repro}
  \]

  \textbf{Stage 1 --- Ingest.} Pulls ETTh1 from a public HTTPS mirror, validates the
  column schema and row count, and writes the raw CSV.

  \textbf{Stage 2 --- Preprocess.} Performs the chronological train/val/test split
  and applies StandardScaler to the OT column, fitting on training data only. The
  scaler is not persisted; since all downstream evaluation is in standardised units,
  no inverse transform is needed.

  \textbf{Stage 3 --- Featurize.} Builds the 45-dimensional feature matrix described
  in Section~IV-B for each partition. Rows with leading NaN values from the longest
  lag window (168 hours) are dropped.

  \textbf{Stage 4 --- Train.} Executes the two-phase hybrid training procedure,
  logs all artefacts to MLflow, and writes a summary to
  \texttt{metrics/metrics.json}.

  \subsection{Experiment Tracking with MLflow}

  Each Stage 4 run opens a parent MLflow run (named \texttt{hybrid\_training}) with
  five nested child runs, one per model. The following are logged per run:
  \begin{itemize}
    \item \textbf{Parameters}: model hyperparameters and pipeline configuration
      (split ratios, lag periods, EWM spans).

    \item \textbf{Metrics}: MAE, MSE, RMSE, MAPE, SMAPE, and $\text{R}^{2}$ for
      Ridge, LightGBM, XGBoost, Ridge+CatBoost, Ridge+ExtraTrees, and the full
      hybrid.

    \item \textbf{Architecture metadata}: base model type, residual model names,
      and feature window sizes, kept alongside metrics for full context.

    \item \textbf{Artefacts}: serialised model files (\texttt{.pkl}) and
      \texttt{metrics/predictions.csv} with per-sample predictions.

    \item \textbf{Registry entry}: the hybrid pipeline is registered and versioned
      in the MLflow Model Registry, enabling controlled promotion from
      \emph{Staging} to \emph{Production}.
  \end{itemize}

  Pairing DVC (pipeline and data state) with MLflow (metric history) yields a
  two-layer audit trail. Given any past prediction one can trace back to the exact
  \texttt{params.yaml} via Git, the exact data split via \texttt{dvc.lock}, and the
  per-model metrics via the associated MLflow run.

  \subsection{Containerisation with Docker}

  Two Docker Compose services are defined. The \textbf{API container} is built on
  \texttt{python:3.12-slim}, copies the trained model binaries at image build time,
  installs Python dependencies, and launches the Uvicorn ASGI server. The
  \textbf{MLflow container} starts MLflow at runtime and mounts host directories
  for the SQLite backend and artefact store. Port 8000 serves predictions; port 5000
  serves the MLflow experiment UI.

  A \texttt{.dockerignore} excludes the virtual environment, Git objects, DVC cache,
  and run artefacts, keeping the image small. Architecturally, the API container is
  stateless (rebuilt from code), while the MLflow container is stateful (data
  persisted via mounted volumes)---a conventional microservice split \cite{docker2013}.

  \subsection{REST API with FastAPI}

  Three endpoints are exposed:
  \begin{itemize}
    \item \texttt{GET /health} --- returns model-load status and API version.

    \item \texttt{GET /metrics} --- reads and returns \texttt{metrics/metrics.json}.

    \item \texttt{POST /predict} --- accepts a start datetime (ISO-8601), a horizon
      \texttt{steps} $\in [1, 720]$, and an optional \texttt{include\_components}
      flag; returns the hybrid forecast and, if requested, per-model breakdowns.
  \end{itemize}

  A \texttt{Predictor} singleton loads all five model files once at startup and keeps
  them in memory, minimising per-request latency. The \texttt{include\_components}
  flag is useful operationally: it lets practitioners inspect the Ridge base
  prediction and individual residual corrections on demand, making anomalous
  predictions easier to diagnose.

  \subsection{CI/CD with GitHub Actions}

  Every push to \texttt{main} triggers three jobs in sequence:
  \begin{enumerate}
    \item \textbf{Lint}: Black, isort, and flake8 run against all Python sources;
      a shared \texttt{setup.cfg} keeps CI and local settings identical.

    \item \textbf{Unit tests}: pytest verifies the health, predict, and metrics
      endpoints; model loading lifecycle; and preprocessing logic.

    \item \textbf{Docker build}: the image is built and a retry-based health-check
      confirms that the container starts and responds correctly.
  \end{enumerate}

  % ── 6. Results ──────────────────────────────────────────────────────────
  \section{Results}

  \subsection{Hybrid Model Performance}

  Table~\ref{tab:hybrid} summarises test-set performance across all trained models.
  Ridge alone proves surprisingly strong---RMSE\,$=0.0765$, $\text{R}^{2}=0.9638$---and
  both LightGBM and XGBoost trained directly on the raw target fall noticeably behind
  ($\text{R}^{2}\approx 0.957$--$0.959$). This gap suggests that the dominant structure
  in ETTh1 is largely linear, and that tree ensembles are handicapped because they
  cannot extrapolate outside their training range.

  Adding CatBoost and ExtraTrees as residual correctors narrows the gap further.
  Ridge\,+\,CatBoost reaches $\text{R}^{2}=0.9639$; Ridge\,+\,ExtraTrees reaches
  $0.9642$. The \textbf{Hybrid Residual} model---averaging both residual
  predictions---achieves RMSE\,$=0.0757$, MAE\,$=0.0520$, SMAPE\,$=7.80\%$, and
  $\text{R}^{2}=0.9645$: the best result on every metric.

  \begin{table}[h]
    \centering
    \caption{Test set performance on ETTh1 (standardised OT)}
    \label{tab:hybrid} \scriptsize
    \setlength{\tabcolsep}{3pt}
    \begin{tabular}{lrrrrr}
      \toprule \textbf{Model Architecture} & \textbf{RMSE}  & \textbf{MAE}   & \textbf{MAPE\,(\%)} & \textbf{SMAPE\,(\%)} & \textbf{R$^{2}$}            \\
      \midrule Ridge (Base Trend)          & 0.0765         & 0.0526         & 10.55               & 7.87                 & \phantom{$-$}0.9638         \\
      LightGBM (Observability)             & 0.0816         & 0.0581         & 11.19               & 8.44                 & \phantom{$-$}0.9588         \\
      XGBoost (Observability)              & 0.0837         & 0.0595         & 11.30               & 8.58                 & \phantom{$-$}0.9567         \\
      \midrule Ridge + CatBoost            & 0.0764         & 0.0525         & \textbf{10.57}      & 7.85                 & \phantom{$-$}0.9639         \\
      Ridge + ExtraTrees                   & 0.0761         & 0.0524         & 11.06               & 7.84                 & \phantom{$-$}0.9642         \\
      \midrule \textbf{Hybrid Residual}    & \textbf{0.0757}& \textbf{0.0520}& 10.74               & \textbf{7.80}        & \phantom{$-$}\textbf{0.9645}\\
      \bottomrule
    \end{tabular}
  \end{table}

  The table also highlights the MAPE anomaly: Ridge\,+\,CatBoost records the lowest
  MAPE (10.57\%) even though the full hybrid achieves lower RMSE, MAE, and SMAPE.
  This is exactly the instability we warned about---a handful of near-zero true values
  disproportionately inflate individual MAPE terms, making ranking by MAPE unreliable
  on standardised data. SMAPE, which symmetrises both sides of the denominator, gives
  the more consistent and interpretable ordering.

  % ── 7. Discussion ──────────────────────────────────────────────────────────
  \section{Discussion}
  \label{sec:discussion}

  \subsection{Why the Hybrid Architecture Works}

  The $\text{R}^{2}$ gap between the hybrid ($0.9645$) and the direct tree models
  ($\approx 0.957$) is not subtle---it amounts to explaining roughly $20\%$ more
  of the remaining unexplained variance. The underlying reason is straightforward:
  tree ensembles are piecewise-constant estimators that average training label values
  at leaf nodes. They cannot produce values outside the training label range, so any
  global trend that carries test observations beyond the training distribution causes
  systematic bias. Ridge, being a global linear model, does not suffer from this. By
  asking the trees to predict only the residual errors, we remove the trend component
  they struggle with and give them a much easier target---one that is closer to
  mean-zero and locally bounded.

  \subsection{Prophet vs.\ Ridge as a Base Model}

  Early in development we experimented with Prophet as the base trend model. The
  problem is structural: after StandardScaling, OT has zero mean and unit variance.
  Prophet's additive formulation interprets this as a signal with near-zero amplitude
  and negligible seasonal components, so it converges to predictions clustered near
  the global mean. The resulting residuals are almost as large as the original signal,
  defeating the purpose of a two-phase architecture. Ridge sidesteps this entirely
  because it fits a linear combination of the already-standardised feature columns and
  never tries to infer trend amplitude from the raw target scale.

  \subsection{Why Average Rather Than Stack}

  We deliberately chose unweighted averaging over stacking for the residual correction
  step. CatBoost (ordered boosting) and ExtraTrees (random-threshold variance reduction)
  are algorithmically orthogonal: they make different approximation errors on the same
  residuals. A simple mean exploits that partial decorrelation without requiring a
  held-out meta-training set or introducing a risk of overfitting at the combination
  layer. The cost is that we cannot learn asymmetric weights that might give more
  credit to whichever corrector is stronger on a given sub-region. Given the marginal
  performance difference between the two individual correctors ($\text{R}^{2}$: $0.9639$
  vs.\ $0.9642$), the simpler approach is defensible.

  \subsection{Reproducibility and Auditability}

  Git, DVC, and MLflow together form a three-layer audit trail. Tracing any historical
  prediction back to its provenance requires only three lookups: the Git commit hash
  for the exact code; the \texttt{dvc.lock} entry for the exact data and parameter
  state; and the MLflow run ID for the per-model metrics. Nothing lives only in a
  notebook or a researcher's local environment---every quantity is persisted and
  queryable.

  \subsection{Limitations}

  Several important limitations bound the scope of these results. First, we evaluate
  only a single ETTh1 dataset at a one-step-ahead horizon; whether the residual
  decomposition benefit holds on shorter series, noisier targets, or multi-step
  horizons is an open question. Second, hyperparameters were chosen from
  domain-informed defaults rather than an optimised search, so the individual model
  numbers should be treated as reasonable baselines rather than upper bounds. Third,
  the Docker deployment lacks horizontal autoscaling, which a high-throughput production
  service would need. Finally, the Prophet limitation is addressable---training on
  unscaled targets before the StandardScaling step would let it function as intended.

  % ── 8. Conclusion ──────────────────────────────────────────────────────────
  \section{Conclusion}

  We built and evaluated a full MLOps-driven Hybrid Residual forecasting system. The
  core technical idea---using Ridge regression to capture the dominant linear trend and
  then training CatBoost and ExtraTrees only on Ridge's errors---yields
  RMSE\,$=0.0757$, SMAPE\,$=7.80\%$, and $\text{R}^{2}=0.9645$ on the ETTh1 test
  set, outperforming every individual model we tested. The gain is not an artefact of
  a complex weighting scheme; it follows from a principled separation of what each
  model family does well.

  Beyond the modelling result, the system itself is a contribution. A four-stage DVC
  pipeline handles data versioning and selective re-execution. MLflow nested runs
  track every model's metrics alongside the parent hybrid experiment. Docker Compose
  packages the serving layer, and a GitHub Actions workflow enforces quality gates on
  every push. The whole system reproduces from raw data to live HTTP inference with
  \texttt{dvc repro}---a single command.

  A secondary methodological finding warrants emphasis: on standardised targets where
  true values frequently approach zero, MAPE is an unreliable ranking metric. We
  recommend reporting SMAPE (or scale-free alternatives) as the primary error measure
  in any forecasting study that involves normalised or zero-mean targets.

  Taken together, these results demonstrate that the model-design and deployment
  concerns of time series forecasting can and should be addressed as an integrated
  engineering problem, not as sequential afterthoughts.

  % ── References ────────────────────────────────────────────────────────────────
  \bibliographystyle{IEEEtran}

  \begin{thebibliography}{99}
    \bibitem{box2015time} G.\ E.\ P.\ Box, G.\ M.\ Jenkins, G.\ C.\ Reinsel, and G.\ M.\ Ljung,
      \textit{Time Series Analysis: Forecasting and Control}, 5th ed. Hoboken,
      NJ: Wiley, 2015.

    \bibitem{taylor2018prophet} S.\ J.\ Taylor and B.\ Letham, ``Forecasting at scale,''
      \textit{The American Statistician}, vol.\ 72, no.\ 1, pp.\ 37--45, 2018.

    \bibitem{ke2017lightgbm} G.\ Ke, Q.\ Meng, T.\ Finley, T.\ Wang, W.\ Chen, W.\ Ma,
      Q.\ Ye, and T.-Y.\ Liu, ``LightGBM: a highly efficient gradient boosting decision
      tree,'' in \textit{Advances in Neural Information Processing Systems (NeurIPS)},
      vol.\ 30, 2017.

    \bibitem{chen2016xgboost} T.\ Chen and C.\ Guestrin, ``XGBoost: a scalable tree
      boosting system,'' in \textit{Proc.\ ACM KDD}, pp.\ 785--794, 2016.

    \bibitem{zhou2021informer} H.\ Zhou, S.\ Zhang, J.\ Peng, S.\ Zhang, J.\ Li, H.\ Xiong,
      and W.\ Zhang, ``Informer: beyond efficient transformer for long sequence
      time-series forecasting,'' in \textit{Proc.\ AAAI}, vol.\ 35, pp.\ 11106--11115,
      2021.

    \bibitem{makridakis2018m4} S.\ Makridakis, E.\ Spiliotis, and V.\ Assimakopoulos,
      ``The M4 competition: results, findings, conclusion and way forward,'' \textit{International
      Journal of Forecasting}, vol.\ 34, no.\ 4, pp.\ 802--808, 2018.

    \bibitem{makridakis2022m5} S.\ Makridakis, E.\ Spiliotis, and V.\ Assimakopoulos,
      ``M5 accuracy competition: results, findings, and conclusions,'' \textit{International
      Journal of Forecasting}, vol.\ 38, no.\ 4, pp.\ 1346--1364, 2022, doi:
      10.1016/j.ijforecast.2021.11.013.

    \bibitem{januschowski2022forecasting} T.\ Januschowski, Y.\ Wang, K.\ Torkkola,
      T.\ Erkkil\"{a}, H.\ Hasson, and J.\ Gasthaus, ``Forecasting with trees,''
      \textit{International Journal of Forecasting}, vol.\ 38, no.\ 4, pp.\ 1473--1481,
      2022, doi: 10.1016/j.ijforecast.2021.10.004.

    \bibitem{lainder2022forecasting} A.\ D.\ Lainder and R.\ D.\ Wolfinger, ``Forecasting
      with gradient boosted trees: augmentation, tuning, and cross-validation strategies,''
      \textit{International Journal of Forecasting}, vol.\ 38, no.\ 4, pp.\ 1426--1433,
      2022, doi: 10.1016/j.ijforecast.2021.12.003.

    \bibitem{zeng2023transformers} A.\ Zeng, M.\ Chen, L.\ Zhang, and Q.\ Xu, ``Are
      transformers effective for time series forecasting?'' in \textit{Proc.\ AAAI
      Conf.\ Artificial Intelligence}, vol.\ 37, no.\ 9, pp.\ 11121--11128, 2023,
      doi: 10.1609/aaai.v37i9.26317.

    \bibitem{nie2023patchtst} Y.\ Nie, N.\ H.\ Nguyen, P.\ Sinthong, and J.\ Kalagnanam,
      ``A time series is worth 64 words: long-term forecasting with transformers,''
      in \textit{Proc.\ International Conference on Learning Representations (ICLR)},
      2023.

    \bibitem{ansari2024chronos} A.\ F.\ Ansari et al., ``Chronos: learning the
      language of time series,'' \textit{Transactions on Machine Learning
      Research}, 2024.

    \bibitem{sculley2015hidden} D.\ Sculley et al., ``Hidden technical debt in
      machine learning systems,'' in \textit{Advances in Neural Information
      Processing Systems (NeurIPS)}, vol.\ 28, 2015.

    \bibitem{kreuzberger2023mlops} D.\ Kreuzberger, N.\ K\"{u}hl, and S.\ Hirschl,
      ``Machine learning operations (MLOps): overview, definition, and architecture,''
      \textit{IEEE Access}, vol.\ 11, pp.\ 31866--31879, 2023.

    \bibitem{berberi2025mlops} L.\ Berberi et al., ``Machine learning operations
      landscape: platforms and tools,'' \textit{Artificial Intelligence Review},
      vol.\ 58, art.\ no.\ 167, 2025, doi: 10.1007/s10462-025-11164-3.

    \bibitem{dvc2020} I.\ Kuprieiev et al., ``DVC: data version control — git for
      data and models,'' \textit{Zenodo}, 2020. [Online]. Available:
      \url{https://dvc.org}

    \bibitem{mlflow2018} M.\ Zaharia et al., ``Accelerating the machine learning
      lifecycle with MLflow,'' \textit{IEEE Data Eng.\ Bull.}, vol.\ 41, no.\ 4, pp.\ 39--45,
      2018.

    \bibitem{docker2013} D.\ Merkel, ``Docker: lightweight Linux containers for
      consistent development and deployment,'' \textit{Linux Journal}, vol.\ 2014,
      no.\ 239, p.\ 2, 2014.

    \bibitem{fastapi2018} S.\ Ramirez, ``FastAPI,'' 2018. [Online]. Available:
      \url{https://fastapi.tiangolo.com}

    \bibitem{bates1969combination} J.\ M.\ Bates and C.\ W.\ J.\ Granger, ``The combination
      of forecasts,'' \textit{Journal of the Operational Research Society}, vol.\ 20,
      no.\ 4, pp.\ 451--468, 1969.

    \bibitem{wolpert1992stacked} D.\ H.\ Wolpert, ``Stacked generalization,'' \textit{Neural
      Networks}, vol.\ 5, no.\ 2, pp.\ 241--259, 1992.

    \bibitem{vogklis2022blending} K.\ Vogklis and I.\ Nasios, ``Blending
      gradient boosted trees and neural networks for point and probabilistic
      forecasting of hierarchical time series,'' \textit{International Journal
      of Forecasting}, vol.\ 38, no.\ 4, pp.\ 1448--1459, 2022, doi: 10.1016/j.ijforecast.2022.01.001.

    \bibitem{hewamalage2021recurrent} H.\ Hewamalage, C.\ Bergmeir, and K.\ Bandara,
      ``Recurrent neural networks for time series forecasting: current status
      and future directions,'' \textit{International Journal of Forecasting}, vol.\ 37,
      no.\ 1, pp.\ 388--427, 2021.

    \bibitem{stan2017} B.\ Carpenter et al., ``Stan: a probabilistic programming language,''
      \textit{Journal of Statistical Software}, vol.\ 76, no.\ 1, pp.\ 1--32, 2017.

    \bibitem{hyndman2018forecasting} R.\ J.\ Hyndman and G.\ Athanasopoulos, \textit{Forecasting:
      Principles and Practice}, 2nd ed. Melbourne: OTexts, 2018. [Online]. Available:
      \url{https://otexts.com/fpp2}
  \end{thebibliography}
\end{document}